% BEGIN EXACT CATALOG SOURCE: thm_9_4
\begin{thmbox}{9.4}
\[
\lvert e^{ix}-1\rvert\leq \lvert x\rvert
\]
for all $x\in\mathbb{R}$.
\end{thmbox}

\textit{Proof}
For fixed $x>0$, we have
\[
\int_0^x e^{iu}\,du
=\left[\frac{e^{iu}}{i}\right]_0^x
=\frac{e^{ix}-1}{i}.
\]
Hence, by the triangle inequality,
\[
\lvert e^{ix}-1\rvert\leq \int_0^x \lvert e^{iu}\rvert\,du=x.
\]
For negative $x$, we can write
\[
\lvert e^{ix}-1\rvert
=\lvert e^{ix}\rvert\lvert 1-e^{i(-x)}\rvert
=\lvert e^{i(-x)}-1\rvert
\leq -x.
\]
Therefore $\lvert e^{ix}-1\rvert\leq \lvert x\rvert$ for all $x\in\mathbb{R}$.
\hfill $\square$
% END EXACT CATALOG SOURCE: thm_9_4

% BEGIN EXACT CATALOG SOURCE: thm_9_5
\begin{thmbox}{9.5}
\textbf{Inversion Formula.}
Let $X$ be a random variable on a probability space $(\Omega,\mathcal{F},P)$, and let $\phi_X(t)$ be the characteristic function of $X$. Denote the push-forward measure of $P$ by $\mu$. For $a<b$,
\[
\mu((a,b))+\frac{\mu(\{a\})}{2}+\frac{\mu(\{b\})}{2}
=\frac{1}{2\pi}\lim_{T\to\infty}
\int_{-T}^{T}\frac{e^{-ita}-e^{-itb}}{it}\phi_X(t)\,dt.
\]
\end{thmbox}

On the left-hand side of the equation in Theorem 9.5, $(a,b)$ denotes an open interval, and $\mu(\{a\})$ and $\mu(\{b\})$ are probabilities at the points $a$ and $b$, respectively. This formulation is necessary to account for possible discontinuities in the cumulative distribution function. If the cdf is continuous, then both $\mu(\{a\})$ and $\mu(\{b\})$ are zero, and $\mu((a,b))=\mu([a,b])$. The limit on the right-hand side is the Cauchy principal value of the integral, and its existence is part of the theorem.

In the proof we use the Dirichlet integral
\[
\lim_{T\to\infty}\int_0^T \frac{\sin u}{u}\,du=\frac{\pi}{2}.
\]
This limit is obtained from the Riemann integral of the function $(\sin u)/u$.

\textit{Proof}
For fixed $T$, consider
\[
I_T \coloneqq
\int_{-T}^{T}\frac{e^{-ita}-e^{-itb}}{it}\phi_X(t)\,dt
=\int_{-T}^{T}\int_{\mathbb{R}}
\frac{e^{-ita}-e^{-itb}}{it}e^{itx}\,d\mu(x)\,dt.
\]
Using Theorem 9.4, we can bound the integrand by
\[
\left\lvert
\frac{e^{-ita}-e^{-itb}}{it}e^{itx}
\right\rvert
=\left\lvert \frac{1}{t}(e^{it(b-a)}-1)\right\rvert
\leq b-a.
\]
Therefore, by Fubini theorem (Theorem 8.5),
\[
I_T=\int_{\mathbb{R}}\int_{-T}^{T}
\frac{e^{-ita}-e^{-itb}}{it}e^{itx}\,dt\,d\mu(x).
\tag{9.1}
\]

Let $f(x,T)$ denote the inner integral in (9.1). The imaginary part of the integrand is odd and the real part is even, so integration over $[-T,T]$ gives
\[
f(x,T)
=2\int_0^T \frac{1}{t}\sin(t(b-x))\,dt
-2\int_0^T \frac{1}{t}\sin(t(a-x))\,dt.
\]
Changing variables gives
\[
f(x,T)
=2\int_{(a-x)T}^{(b-x)T}\frac{\sin u}{u}\,du.
\]
There are four cases. If $a<x<b$, then $f(x,T)\to 2\pi$. If $x=a$ or $x=b$, then $f(x,T)\to \pi$. If $x<a$ or $x>b$, then $f(x,T)\to 0$. Thus
\[
\lim_{T\to\infty} f(x,T)=
\begin{cases}
2\pi, & x\in(a,b),\\
\pi, & x=a\text{ or }x=b,\\
0, & \text{otherwise}.
\end{cases}
\]
The integral $\int_0^v (\sin u)/u\,du$ converges as $v\to\infty$ and is continuous as a function of $v$. Hence $\lvert f(x,T)\rvert$ is bounded by a constant independent of $x$ and $T$. Applying the complex version of the dominated convergence theorem (Theorem 7.6), we get
\[
\lim_{T\to\infty} I_T
=2\pi\int 1_{(a,b)}(x)\,d\mu(x)
+\pi\int 1_{\{a\}}(x)\,d\mu(x)
+\pi\int 1_{\{b\}}(x)\,d\mu(x).
\]
Dividing by $2\pi$ gives the inversion formula.
\hfill $\square$
% END EXACT CATALOG SOURCE: thm_9_5

% BEGIN EXACT CATALOG SOURCE: thm_9_6
Using the inversion formula, we can show that the characteristic function of a random variable uniquely determines the cumulative distribution function.

\begin{thmbox}{9.6}
\textbf{Uniqueness Theorem.}
Let $X$ and $Y$ be random variables with characteristic functions $\phi_X(t)$ and $\phi_Y(t)$, respectively. If $\phi_X(t)=\phi_Y(t)$ for all $t\in\mathbb{R}$, then $X$ and $Y$ have the same distribution.
\end{thmbox}

\textit{Proof}
Let $F_X(x)$ and $F_Y(y)$ denote the cumulative distribution functions of $X$ and $Y$, respectively. If $F_X$ and $F_Y$ are both continuous at $a$ and $b$, the inversion formula can be simplified to
\[
P(a<X<b)
=\frac{1}{2\pi}\lim_{T\to\infty}\int_{-T}^{T}
\frac{e^{-iat}-e^{-ibt}}{it}\phi_X(t)\,dt,
\]
and
\[
P(a<Y<b)
=\frac{1}{2\pi}\lim_{T\to\infty}\int_{-T}^{T}
\frac{e^{-iat}-e^{-ibt}}{it}\phi_Y(t)\,dt.
\]
Since $\phi_X(t)=\phi_Y(t)$ for all $t$, we have $P(a<X<b)=P(a<Y<b)$.

Because there are at most countably many discontinuity points in the cumulative distribution function of a random variable (see Exercise 3.5), the Lebesgue--Stieltjes measures induced by $F_X$ and $F_Y$ agree on a collection of open intervals $(a,b)$, where $a$ and $b$ range over a dense set in $\mathbb{R}$. This collection of intervals generates the Borel algebra on $\mathbb{R}$. Therefore, the two Lebesgue--Stieltjes measures are the same. Consequently, the cdf's of $X$ and $Y$ are the same (see Theorem 3.9).
\hfill $\square$
% END EXACT CATALOG SOURCE: thm_9_6

% BEGIN EXACT CATALOG SOURCE: prob_9_6
\textbf{Problem 9.6} Let $X_1,\ldots,X_n$ be i.i.d. Cauchy random variables. Show that $(X_1+\cdots+X_n)/n$ is Cauchy with the same distribution as $X_1$.
% END EXACT CATALOG SOURCE: prob_9_6

% BEGIN EXACT CATALOG SOURCE: prob_9_8
\textbf{Problem 9.8} Prove that $\phi_X(2\pi)=1$ if and only if $P(X\in\mathbb{Z})=1$.
% END EXACT CATALOG SOURCE: prob_9_8
